\begin{abstract}
\noindent
An issuer that lends its redemption reserve against its own locked tokens can pay every redemption today and still fail after a borrower repays part of a loan and withdraws collateral. What the issuer may safely lend is therefore set by the withdrawal rights that its loans promise. We characterize coverage after every permitted repayment for a fixed issuance curve. With proportional release and no extra reserve, every individual loan must be bounded by the floor value of its collateral, whatever the allocation of debt. With a committed reserve, a sorted debt-per-collateral envelope gives the exact reserve requirement and the maximum aggregate credit, and equal debt ratios minimize the requirement. Our main result ties that requirement to incentives: under stated zero-fee coalition assumptions, the required reserve equals the most that borrowers can gain by repaying in part, redeeming what is released and abandoning the rest, so that at the maximum admitted credit the strategy can exhaust a reserve that others funded. A contract with nonlinear release lends against the reserve value of a whole lock without that exposure, at the cost of less flexible withdrawals and stronger conditions on repricing and splitting.

These results rest on a safety theorem for the floor-capped contract: for an arbitrary issuance schedule, a three-part invariant gives redemption coverage and loan payability on every path without any repayment, and its preservation under a complete operation set is formalized in Lean~4. For that contract, a price-taking borrower who values cash flows and has no use for the locked tokens in the meantime weakly prefers to defer repayment, since the loan has no carrying cost, and a stationary economy with given arrivals determines utilization and floor growth; its own cash flows cannot support voluntary entry when growth is below the discount rate and the loan limit is below the net floor. Further results separate reserve safety from redistribution and bound sandwich and repricing extraction.

An integer reference model, Solidity tests and exact rational checks support distinct parts of the analysis. Descriptive evidence from \ELogs{} event logs of Baseline's credit system shows that \ELeverAmtPct\% of new gross debt finances token purchases, and a reconciliation of fee distributions separates amounts allocated to staking from rewards claimed. These observations motivate a theory of leveraged demand; they do not identify one.
\end{abstract}
